% Generated draft: requires mathematical and source review.
\documentclass[11pt]{article}
\usepackage[margin=25mm]{geometry}
\usepackage{fontspec}
\setmainfont{Latin Modern Roman}
\usepackage{amsmath,amssymb,mathtools,unicode-math}
\setmathfont{Latin Modern Math}
\usepackage{enumitem,needspace,xurl,hyperref}
\setlength{\parindent}{0pt}
\setlength{\parskip}{5pt}
\setlength{\emergencystretch}{3em}
\newcommand{\N}{\mathbb{N}}
\newcommand{\Z}{\mathbb{Z}}
\newcommand{\Q}{\mathbb{Q}}
\newcommand{\R}{\mathbb{R}}
\newcommand{\C}{\mathbb{C}}
\newcommand{\F}{\mathbb{F}}
\newcommand{\A}{\mathbb{A}}
\newcommand{\PP}{\mathbb{P}}
\newcommand{\E}{\mathbb{E}}
\newcommand{\eps}{\varepsilon}
\newcommand{\dd}{\,\mathrm{d}}
\begin{document}
\section*{UnsolvedMath ID: 30001452}
\textbf{Source problem code:} catalogue entry
\subsection{Definitions and mathematical statement}
Use the following explicit standard conventions to interpret the supplied question. Let \(n\geq 1\) be an integer, let \(F\) be a field, and let \(q\in F^{\times}\). Write \(H_n(F,q)\) for the Iwahori--Hecke algebra of the symmetric group \(\mathfrak S_n\), of Coxeter type \(A_{n-1}\). It is the unital associative \(F\)-algebra generated by \(T_1,\ldots,T_{n-1}\), with relations
\[
(T_i-q)(T_i+1)=0,\qquad
T_iT_{i+1}T_i=T_{i+1}T_iT_{i+1},\qquad
T_iT_j=T_jT_i\quad (|i-j|>1),
\]
for indices for which the expressions are defined; set \(H_1(F,q)=F\). A partition \(\lambda\vdash n\) is a finite weakly decreasing sequence of positive integers with sum \(n\). Let \(S_F^{\lambda}(q)\) denote the standard Specht module of \(H_n(F,q)\) indexed by \(\lambda\). A module is irreducible if it is nonzero and has no nonzero proper submodule.

\textbf{Classification problem.} For every \(n\geq 1\), classify all triples
\[
(\lambda,F,q),\qquad \lambda\vdash n,\quad F\text{ a field},\quad q\in F^{\times},
\]
for which \(S_F^{\lambda}(q)\) is irreducible over \(F\).

For the terminology in the source, put \(p=\operatorname{char}F\), where \(p=0\) or \(p\) is prime, and define the quantum characteristic by
\[
e=\min\{r\geq 2:1+q+\cdots+q^{r-1}=0\text{ in }F\},
\]
with \(e=\infty\) if this set is empty. The supplied record describes the classification equivalently as that of \((e,p)\)-irreducible partitions. Here this phrase refers to partitions whose corresponding Specht modules are irreducible for the parameters under consideration. The claim that this property depends only on \((e,p)\), rather than on further details of \(F\) and \(q\), is part of the supplied equivalence; no justification or precise range of that equivalence is included in the record.

\textbf{Unresolved qualifications.} The presentation and the restriction \(q\ne 0\) above make explicit a usual convention for Iwahori--Hecke algebras; the supplied statement does not specify a presentation, parameter normalization, or whether a zero-parameter specialization is intended. If \(q=0\) is intended, its inclusion requires a separate clarification. No algebraic-closure or splitting-field hypothesis is supplied, and the stated target is irreducibility over \(F\). The supplied background also refers to a general Iwahori--Hecke extension and to type \(B\), whereas the extracted mathematical question explicitly concerns type \(A\). It does not identify which remaining cases are intended as unresolved.
\subsection{Short English statement}
Classify the partitions, fields, and Hecke parameters for which a type-\(A\) Specht module is irreducible. The formulation uses explicit standard nonzero-parameter conventions; the intended remaining cases and the scope of the source's equivalent \((e,p)\) terminology require review.
\subsection{Sources}
\begin{itemize}
\item Supplied local UnsolvedMath registry record, permanent ID \(30001452\), title \emph{Irreducible Specht Modules for Type-A Hecke Algebras}, source problem code \texttt{catalogue entry}, in the Oberwolfach Reports Open Problems collection. The record gives proposed year \(2010\) and supplies no individual proposer.
\item Original report: \emph{Combinatorial Representation Theory} (2010), Oberwolfach Reports, pp.\ 799--882, DOI \texttt{10.4171/owr/2010/15}. Supplied source and literature link: \url{https://doi.org/10.4171/owr/2010/15}.
\item Supplied literature reference: Fayers, irreducible Specht modules in type A, \url{https://doi.org/10.1016/j.aim.2004.06.001}.
\item Supplied literature reference: Fayers, type B, \url{https://arxiv.org/abs/1005.5660}.
\item Supplied dated literature curation, checked \texttt{2026-08-21}: the record labels the problem partially solved and assesses that the symmetric-group Specht-module problem is solved, substantial type-\(A\) and type-\(B\) Iwahori--Hecke classifications are known, and the extracted general Iwahori--Hecke extension is not completely classified. This assessment is supplied dated curation derived from the OWR material, not independently verified here; its time-sensitive status and cited claims require independent verification.
\item The registry labels the statement \texttt{corrected\_verified} and reports that its mathematical content was checked against the official Oberwolfach report, with notation, context, and question form normalized without changing the source-level claim. This is supplied verification metadata, not an independent source check for this draft.
\end{itemize}
\end{document}
